% Generated draft: requires mathematical and source review.
\documentclass[11pt]{article}
\usepackage[margin=25mm]{geometry}
\usepackage{fontspec}
\setmainfont{Latin Modern Roman}
\usepackage{amsmath,amssymb,mathtools,unicode-math}
\setmathfont{Latin Modern Math}
\usepackage{enumitem,needspace,xurl,hyperref}
\setlength{\parindent}{0pt}
\setlength{\parskip}{5pt}
\setlength{\emergencystretch}{3em}
\newcommand{\N}{\mathbb{N}}
\newcommand{\Z}{\mathbb{Z}}
\newcommand{\Q}{\mathbb{Q}}
\newcommand{\R}{\mathbb{R}}
\newcommand{\C}{\mathbb{C}}
\newcommand{\F}{\mathbb{F}}
\newcommand{\A}{\mathbb{A}}
\newcommand{\PP}{\mathbb{P}}
\newcommand{\E}{\mathbb{E}}
\newcommand{\eps}{\varepsilon}
\newcommand{\dd}{\,\mathrm{d}}
\begin{document}
\section*{UnsolvedMath ID: 30002398}
\textbf{Source problem code:} catalogue entry
\subsection{Definitions and mathematical statement}
Let \(A\) be a \(C^*\)-algebra, with neither simplicity nor unitality assumed. No separability hypothesis is explicitly supplied.

A completely positive map is called order zero if it preserves orthogonality of positive elements. For \(n\in\{0,1,2,\ldots\}\), write \(\dim_{\mathrm{nuc}}(A)\leq n\) if, for every finite subset \(\mathcal F\subseteq A\) and every \(\varepsilon>0\), there exist finite-dimensional \(C^*\)-algebras \(D_0,\ldots,D_n\), a completely positive contractive map \(\psi:A\to\bigoplus_{j=0}^n D_j\), and completely positive contractive order-zero maps \(\phi_j:D_j\to A\) such that, writing \(\psi_j\) for the components of \(\psi\),
\[
\left\|a-\sum_{j=0}^n\phi_j\psi_j(a)\right\|<\varepsilon
\qquad(a\in\mathcal F).
\]
Finite nuclear dimension means that this holds for some finite \(n\).

Here amenability is understood in the standard \(C^*\)-algebra sense, equivalently nuclearity: for every finite \(\mathcal F\subseteq A\) and \(\varepsilon>0\), there are a finite-dimensional \(C^*\)-algebra \(D\) and completely positive contractive maps \(\alpha:A\to D\) and \(\beta:D\to A\) satisfying \(\|\beta\alpha(a)-a\|<\varepsilon\) for all \(a\in\mathcal F\). Let \(\mathcal Z\) denote the Jiang--Su algebra. The algebra \(A\) is \(\mathcal Z\)-stable if
\[
A\cong A\otimes_{\min}\mathcal Z.
\]

\textbf{Source hypothesis requiring clarification.} Denote by \(\mathsf H(A)\) the source condition that ``no ideal has a type-I representation.'' Ideals are understood as closed two-sided ideals. The supplied record does not define a type-I representation or state the nonzero conventions for the ideal and representation. In particular, it does not specify whether type I concerns the image \(\pi(I)\) of a representation \(\pi:I\to B(H)\), or the generated von Neumann algebra \(\pi(I)''\). These interpretations must not be silently identified. Zero representations must also be addressed to obtain an intended nontrivial hypothesis. Accordingly, \(\mathsf H\) retains the source wording rather than asserting a uniquely determined representation-theoretic condition.

\textbf{Question, conditional on resolving the source hypothesis.} Is it true that, for every \(C^*\)-algebra \(A\) satisfying \(\mathsf H(A)\),
\[
\dim_{\mathrm{nuc}}(A)<\infty
\quad\Longleftrightarrow\quad
\bigl(A\text{ is amenable and }A\cong A\otimes_{\min}\mathcal Z\bigr),
\]
without assuming simplicity or unitality? The quantification here adds no separability assumption; whether the original workshop context implicitly imposed one remains unresolved. This is a statement-only draft, and its intended class of algebras requires human mathematical review.
\subsection{Short English statement}
The question asks whether finite nuclear dimension is equivalent to amenability together with absorption of the Jiang--Su algebra for the source's class of \(C^*\)-algebras, allowing nonsimple and nonunital algebras. The condition excluding type-I representations of ideals needs clarification, and no separability assumption is supplied. The registry labels the problem open; this draft does not independently verify that status.
\subsection{Sources}
\begin{itemize}
\item Local UnsolvedMath registry, supplied complete record for permanent ID \(30002398\), title \emph{Nuclear Dimension and Z-Stability Beyond Simple Algebras}, source problem code \texttt{catalogue entry}, in the Oberwolfach Reports Open Problems collection. The supplied proposed year is 2013; no individual proposer is identified.
\item Supplied source citation: \emph{C*-Algebren (2014)}, Oberwolfach Reports, DOI \texttt{10.4171/owr/2013/43}, pp.\ 2423--2500. \url{https://doi.org/10.4171/owr/2013/43}. The background attributes the collection to workshop open-problem discussions documented in Oberwolfach Reports. The citation's 2014 label and its 2013 DOI/year metadata are retained without reconciliation.
\item Supplied literature source, ``Oberwolfach Report'': \url{https://ems.press/content/serial-article-files/46475}.
\item Supplied literature source, ``Robert--Tikuisis'': \url{https://arxiv.org/abs/1308.2941}.
\item Supplied literature source, ``Enstad--Vilalta'': \url{https://arxiv.org/abs/2503.18088}.
\item Supplied literature assessment dated exactly \texttt{2026-08-21}: the general nonsimple Toms--Winter implication remains open, while available results cover important special classes rather than the stated generality. This is supplied dated curation derived from the OWR record, not an independently verified literature assessment; time-sensitive status and cited claims require independent verification.
\item The registry gives statement status \texttt{corrected\_verified} and reports that the mathematical content was checked against the official Oberwolfach report, with notation, context, and question form normalized without changing the source-level claim. These are supplied provenance assertions; no source was visited or independently checked for this draft.
\end{itemize}
\end{document}
